\begin{lemma}
Let \(I\subseteq [n]\), put \(k=|I|\), and write
\[
\ell_R=|\pi_R(I)|,\qquad \ell_B=|\pi_B(I)|.
\]
Assume \(\varepsilon_1\ge0\),
\(\noderef{ParameterHierarchyEventualComparisons}(\varepsilon,\varepsilon_1,
\varepsilon_2,n_0)\), \(n>n_0\), and
\[
k=\noderef{TwoBiteNaturalIndependenceScale}(1+\varepsilon,n).
\]
Suppose also that the sample parameters are the rounded hierarchy parameters
\[
m=\noderef{TwoBiteNaturalReducedVertexCount}(n),
\qquad
p=\noderef{TwoBiteEdgeProbability}\!\left(\frac12,n\right).
\]
If \(\noderef{FiberAndDegreeEvent}\) holds and the medium and small closed-pair
events \(\noderef{TwoBiteMediumClosedPairsEvent}\) and
\(\noderef{TwoBiteSmallClosedPairsEvent}\) hold at the same parameters, then
after the staged revelation procedure at least
\[
f(\ell_R,\ell_B,k,p,n)-10\varepsilon_1k^2
\]
open projected pairs remain unrevealed, where
\(f=\noderef{ProjectionOpenPairFunction}\).  Equivalently,
\[
f(\ell_R,\ell_B,k,p,n)-10\varepsilon_1k^2
\le |\noderef{TwoBiteStagedOpenPairs}(\omega,\varepsilon,I)|.
\]
\end{lemma}
\begin{proof}
All sets in the proof are finite subsets of the same tagged projected-pair
ambient set.  Put
\[
P_I=\noderef{TwoBiteProjectionPairSet}(\omega,I),
\]
and write \(C(e)\) for
\(\noderef{TwoBiteProjectionPairClosed}(\omega,I,e)\),
\(C_{\mathrm{sc}}(e)\) for
\(\noderef{TwoBiteProjectionPairSameColorClosed}(\omega,I,e)\), \(R(e)\)
for \(\noderef{TwoBiteProjectionPairTouched}(\omega,\varepsilon,I,e)\), and
\(Q(e)\) for
\(e\in\noderef{TwoBitePreTerminalRecordedPairs}(\omega,\varepsilon,I)\).  With
this notation the definitions give
\[
\noderef{TwoBiteUnclosedProjectedPairs}(\omega,I)
  =\{e\in P_I:\neg C(e)\},
\]
\[
\noderef{TwoBiteTouchedProjectedPairs}(\omega,\varepsilon,I)
  =\{e\in P_I:R(e)\},
\]
and
\[
\noderef{TwoBiteStagedOpenPairs}(\omega,\varepsilon,I)
  =\{e\in P_I:\neg C_{\mathrm{sc}}(e)\wedge\neg R(e)\wedge\neg Q(e)\}.
\]
Let
\[
U_I=\noderef{TwoBiteUnclosedProjectedPairs}(\omega,I),\qquad
T_I=\noderef{TwoBiteTouchedProjectedPairs}(\omega,\varepsilon,I).
\]
By \(\noderef{TwoBiteSameColorClosedImpliesClosed}\), applied to this same
sample \(\omega\), this same set \(I\), and an arbitrary tagged coordinate
\(e\), same-color closedness in the one-sided \(C^+(I)\) sense implies full
closedness in the \(X_x(I)\) sense.  In the notation above, the helper gives
the pointwise implication
\[
C_{\mathrm{sc}}(e)\Longrightarrow C(e).
\tag{A0}
\]
By the definition of \(\noderef{TwoBitePreTerminalRecordedPairs}\), a
pre-terminal recorded coordinate has one of two forms.  In the incident
coordinate branch, it is a same-color coordinate with one endpoint
staged-revealed and the other endpoint projection-contained.  Unfolding
\(\noderef{TwoBiteProjectionPairTouched}\) shows that such a coordinate is
touched.  In the \(X_x(I)\)-pair branch, it is represented by a pair in
\(X_x(I)\) for some staged-revealed vertex \(x\);
unfolding \(\noderef{TwoBiteProjectionPairClosed}\) shows that it is closed.
Thus
\[
Q(e)\Longrightarrow C(e)\vee R(e).
\tag{A1}
\]
For a fixed \(e\), the ordinary propositional contrapositive of (A0) is
\[
\neg C(e)\Longrightarrow \neg C_{\mathrm{sc}}(e).
\tag{A0'}
\]
\[
\neg C(e)\wedge\neg R(e)\Longrightarrow \neg Q(e)
\tag{A1'}
\]
is the propositional contrapositive of (A1).
The definitions of \(\noderef{TwoBiteUnclosedProjectedPairs}\),
\(\noderef{TwoBiteTouchedProjectedPairs}\), and
\(\noderef{TwoBiteStagedOpenPairs}\), together with (A0') and (A1'), imply the inclusion
\[
U_I\setminus T_I
\subseteq
\noderef{TwoBiteStagedOpenPairs}(\omega,\varepsilon,I).
\tag{A}
\]
Indeed, if \(e\in U_I\setminus T_I\), then \(e\in U_I\) and \(e\notin T_I\).
The first membership unfolds to \(e\in P_I\) and \(\neg C(e)\).  The
contrapositive (A0') gives \(\neg C_{\mathrm{sc}}(e)\).  To get the touched
exclusion, suppose for contradiction that
\(R(e)\) held.  Since we already know
\(e\in P_I\), the defining filter identity
\(T_I=\{e\in P_I:R(e)\}\) would put \(e\) in \(T_I\), contradicting
\(e\notin T_I\).  Hence \(\neg R(e)\).  Applying (A1') to
\(\neg C(e)\) and \(\neg R(e)\) gives \(\neg Q(e)\).  Thus \(e\in P_I\),
\(\neg C_{\mathrm{sc}}(e)\), \(\neg R(e)\), and \(\neg Q(e)\), which is
exactly the defining membership condition for
\(\noderef{TwoBiteStagedOpenPairs}(\omega,\varepsilon,I)\).

For the finite sets \(U_I\) and \(T_I\), the set \(U_I\) is the disjoint union
of the elements outside \(T_I\) and the elements inside \(T_I\):
\[
U_I=(U_I\setminus T_I)\,\dot\cup\,(U_I\cap T_I).
\]
The two pieces are disjoint because no element can both lie in \(T_I\) and lie
outside \(T_I\), and every element of \(U_I\) lies in exactly one of the two
pieces.  Therefore the natural-cardinality equality is
\[
|U_I|=|U_I\setminus T_I|+|U_I\cap T_I|.
\]
The inclusion \(U_I\cap T_I\subseteq T_I\) gives
\[
|U_I\cap T_I|\le |T_I|.
\]
Adding \(|U_I\setminus T_I|\) to both sides gives
\[
|U_I\setminus T_I|+|U_I\cap T_I|
\le |U_I\setminus T_I|+|T_I|.
\]
Using the preceding equality for the left-hand side,
\[
|U_I|\le |U_I\setminus T_I|+|T_I|.
\]
Cast this natural-cardinality inequality to real numbers.  Since the cast
\(\mathbb N\to\mathbb R\) is order preserving and respects addition, we get
\[
(|U_I|:\mathbb R)
\le (|U_I\setminus T_I|:\mathbb R)+(|T_I|:\mathbb R).
\]
Subtracting \((|T_I|:\mathbb R)\) from both sides gives the real-valued
cardinality inequality used below:
\[
(|U_I|:\mathbb R)-(|T_I|:\mathbb R)
\le (|U_I\setminus T_I|:\mathbb R).
\tag{B}
\]
Combining the inclusion (A) with monotonicity of finite-set cardinality gives
\[
(|U_I\setminus T_I|:\mathbb R)
\le
\left(\left|\noderef{TwoBiteStagedOpenPairs}(\omega,\varepsilon,I)\right|:
\mathbb R\right).
\tag{C}
\]
Therefore, from (B) and (C),
\[
\left(\left|\noderef{TwoBiteStagedOpenPairs}(\omega,\varepsilon,I)\right|:
\mathbb R\right)
\ge (|U_I|:\mathbb R)-(|T_I|:\mathbb R).
\]

Now apply \(\noderef{ProjectionClosedPairLossBound}\) with exactly the same
objects as in the present lemma: the sample is \(\omega\), the vertex set is
\(I\), the parameters are \(\varepsilon,\varepsilon_1,\varepsilon_2,n_0\), the
ambient sample has the same \(n,m,p\), and the rounded identities supplied to
the child are
\[
m=\noderef{TwoBiteNaturalReducedVertexCount}(n),\qquad
p=\noderef{TwoBiteEdgeProbability}\!\left(\frac12,n\right),\qquad
|I|=\noderef{TwoBiteNaturalIndependenceScale}(1+\varepsilon,n).
\]
These are precisely the rounded \(m\), rounded \(p\), and
\(k=|I|\) hypotheses of the active statement.  The remaining hypotheses of the
child are also present verbatim: \(0\le\varepsilon_1\),
\(\noderef{ParameterHierarchyEventualComparisons}
(\varepsilon,\varepsilon_1,\varepsilon_2,n_0)\), \(n>n_0\),
\(\noderef{FiberAndDegreeEvent}(\omega,\varepsilon_2)\), and the medium and
small closed-pair events with the child theorem's explicit parameter
\(k:=|I|\):
\[
\noderef{TwoBiteMediumClosedPairsEvent}(\varepsilon,\varepsilon_1,\omega)
\text{ at } k=|I|,
\qquad
\noderef{TwoBiteSmallClosedPairsEvent}(\varepsilon,\varepsilon_1,\omega)
\text{ at } k=|I|.
\]
The conclusion of the child, with this same \(\omega\) and \(I\), is
\[
f\!\left(((I.image(\pi_R^\omega)).card:\mathbb R),
((I.image(\pi_B^\omega)).card:\mathbb R),
(|I|:\mathbb R),p,(n:\mathbb R)\right)
-9\varepsilon_1(|I|:\mathbb R)^2
\le (|U_I|:\mathbb R).
\]
By the definitions fixed at the start of the statement,
\[
\ell_R=((I.image(\pi_R^\omega)).card:\mathbb R),\qquad
\ell_B=((I.image(\pi_B^\omega)).card:\mathbb R),\qquad
k=(|I|:\mathbb R).
\]
Thus the displayed child conclusion is exactly
\[
f(\ell_R,\ell_B,k,p,n)-9\varepsilon_1k^2\le (|U_I|:\mathbb R).
\tag{D}
\]

Apply \(\noderef{StagedTouchedPairsBound}\) with the same sample \(\omega\),
the same set \(I\), the same parameters
\(\varepsilon,\varepsilon_1,\varepsilon_2,n_0\), and the same scale identity
\[
|I|=\noderef{TwoBiteNaturalIndependenceScale}(1+\varepsilon,n).
\]
This child does not require the rounded \(m\) and \(p\) identities, so those
ambient identities are simply unchanged and unused in this application.  Its
required hypotheses \(0\le\varepsilon_1\), the same eventual-comparisons
package, \(n>n_0\), and
\(\noderef{FiberAndDegreeEvent}(\omega,\varepsilon_2)\) are exactly among the
active assumptions.  Its conclusion is stated for
\(\noderef{TwoBiteTouchedProjectedPairs}(\omega,\varepsilon,I)\), which is
the set \(T_I\) defined above, so it gives
\[
(|T_I|:\mathbb R)\le \varepsilon_1k^2.
\tag{E}
\]

Combine the finite-set inequality with the two child estimates as follows.
From (B) and (C),
\[
\left(\left|\noderef{TwoBiteStagedOpenPairs}(\omega,\varepsilon,I)\right|:
\mathbb R\right)
\ge (|U_I|:\mathbb R)-(|T_I|:\mathbb R).
\tag{F}
\]
Inequality (D) says that \((|U_I|:\mathbb R)\) is at least
\[
f(\ell_R,\ell_B,k,p,n)-9\varepsilon_1k^2,
\]
and (E) says that subtracting \((|T_I|:\mathbb R)\) loses at most
\(\varepsilon_1k^2\).  Therefore
\[
(|U_I|:\mathbb R)-(|T_I|:\mathbb R)
\ge \bigl(f(\ell_R,\ell_B,k,p,n)-9\varepsilon_1k^2\bigr)
    -\varepsilon_1k^2.
\tag{G}
\]
Combining (F) and (G) gives
\[
\left(\left|\noderef{TwoBiteStagedOpenPairs}(\omega,\varepsilon,I)\right|:
\mathbb R\right)
\ge
\bigl(f(\ell_R,\ell_B,k,p,n)-9\varepsilon_1k^2\bigr)-\varepsilon_1k^2.
\]
The right-hand side is
\[
f(\ell_R,\ell_B,k,p,n)-10\varepsilon_1k^2,
\]
because \(9\varepsilon_1k^2+\varepsilon_1k^2=10\varepsilon_1k^2\).  Thus
\[
f(\ell_R,\ell_B,k,p,n)-10\varepsilon_1k^2
\le
\left(\left|\noderef{TwoBiteStagedOpenPairs}(\omega,\varepsilon,I)\right|:
\mathbb R\right).
\]
This is the asserted lower bound.
\end{proof}
